%%%%%%%%%%%%%%%%%%%%%%%%%%%%%%%%%%%%%%%%%%%%%%%%%%%%%%%%%%%%%%%%%%%%%%
% LaTeX Example: Project Report
%
% Source: http://www.howtotex.com
%
% Feel free to distribute this example, but please keep the referral
% to howtotex.com
% Date: March 2011 
% 
%%%%%%%%%%%%%%%%%%%%%%%%%%%%%%%%%%%%%%%%%%%%%%%%%%%%%%%%%%%%%%%%%%%%%%
% How to use writeLaTeX: 
%
% You edit the source code here on the left, and the preview on the
% right shows you the result within a few seconds.
%
% Bookmark this page and share the URL with your co-authors. They can
% edit at the same time!
%
% You can upload figures, bibliographies, custom classes and
% styles using the files menu.
%
% If you're new to LaTeX, the wikibook is a great place to start:
% http://en.wikibooks.org/wiki/LaTeX
%
%%%%%%%%%%%%%%%%%%%%%%%%%%%%%%%%%%%%%%%%%%%%%%%%%%%%%%%%%%%%%%%%%%%%%%
% Edit the title below to update the display in My Documents
%\title{Project Report}
%
%%% Preamble
\documentclass[paper=a4, fontsize=11pt]{scrartcl}
\usepackage[T1]{fontenc}
\usepackage{fourier}

\usepackage[english]{babel}															% English language/hyphenation
\usepackage[protrusion=true,expansion=true]{microtype}	
\usepackage{amsmath,amsfonts,amsthm} % Math packages
\usepackage[pdftex]{graphicx}	
\usepackage{url}
\usepackage{titlesec}
\usepackage{mathtools}
%\usepackage{lstlisting}


%%% Custom headers/footers (fancyhdr package)
\usepackage{fancyhdr}
\pagestyle{fancyplain}
\fancyhead{}											% No page header
\fancyfoot[L]{}											% Empty 
\fancyfoot[C]{}											% Empty
\fancyfoot[R]{\thepage}									% Pagenumbering
\renewcommand{\headrulewidth}{0pt}			% Remove header underlines
\renewcommand{\footrulewidth}{0pt}				% Remove footer underlines
\setlength{\headheight}{13.6pt}
\usepackage[margin=0.5in]{geometry}

%%% Equation and float numbering
\numberwithin{equation}{section}		% Equationnumbering: section.eq#
\numberwithin{figure}{section}			% Figurenumbering: section.fig#
\numberwithin{table}{section}				% Tablenumbering: section.tab#

%% Section Format
\titleformat*{\section}{\LARGE\bfseries}
\titleformat*{\subsection}{\Large\bfseries}
\titleformat*{\subsubsection}{\large\bfseries}
\titleformat*{\paragraph}{\large\bfseries}
\titleformat*{\subparagraph}{\large\bfseries}

\DeclarePairedDelimiter{\ceil}{\lceil}{\rceil}

%%% Maketitle metadata
\newcommand{\horrule}[1]{\rule{\linewidth}{#1}} 	% Horizontal rule

\title{
		%\vspace{-1in} 	
		\usefont{OT1}{bch}{b}{n}
		\normalfont \normalsize \textsc{University of Purdue\\
      Data Communication and Computer Networks\\
      CS 53600 Spring 2020} \\ [25pt]
		\horrule{0.5pt} \\[0.4cm]
		\huge Assignment 4\\
		\horrule{2pt} \\[0.5cm]
}
\author{
		\normalfont\normalsize Pedro da Costa Abreu Júnior\\ [-5pt]
    \normalsize Student ID 0030878383\\[-5pt]
    \normalsize \today
}
\date{}


%%% Begin document
\begin{document}
\maketitle

\section*{Problem 1}

\section*{Problem 2}
\subsection*{a}
Only traffic arriving from hosts h1 and h6 should be delivered to hosts h3 or h4
(i.e., that arriving traffic from hosts h2 and h5 is blocked).

\begin{center}
\begin{tabular}{ | c | c | c | c | c | c | c | c | c | c | c | c | }
  \hline
Switch & Port & MAC src & MAC dst & Eth type & VLAN ID & IP Src & IP Dst & IP Prot & TCP sport & TCP dport & Action \\
  \hline
  s2 & * & * & * & * & * & 10.1.0.2 & * & * & * & * & drop \\
  \hline
  s2 & * & * & * & *  & * & 10.1.0.5 & * & * & * & * & drop \\
  \hline
\end{tabular}
\end{center}


\subsection*{b}
Only TCP traffic is allowed to be delivered to host h3 or host h4 (i.e., that
UDP traffic is blocked).

\begin{center}
\begin{tabular}{ | c | c | c | c | c | c | c | c | c | c | c | c | }
  \hline
Switch & Port & MAC src & MAC dst & Eth type & VLAN ID & IP Src & IP Dst & IP Prot & TCP sport & TCP dport & Action \\
  \hline
  s2 & * & * & * & * & * & * & * & UDP & * & * & drop \\
  \hline
\end{tabular}
\end{center}

\subsection*{c}
Only traffc destined to h3 is to be delivered (i.e., all traffic to h4 is
blocked).

\begin{center}
\begin{tabular}{ | c | c | c | c | c | c | c | c | c | c | c | c | }
  \hline
Switch & Port & MAC src & MAC dst & Eth type & VLAN ID & IP Src & IP Dst & IP Prot & TCP sport & TCP dport & Action \\
  \hline
  s2 & * & * & * & * & * & * & 10.1.0.4 & * & * & * & drop \\
  \hline
\end{tabular}
\end{center}


\subsection*{d}
Only UDP traffic from h1 and destined to h3 is to be delivered. All other traffic is blocked.



\begin{center}
\begin{tabular}{ | c | c | c | c | c | c | c | c | c | c | c | l | }
  \hline
Switch & Port & MAC src & MAC dst & Eth type & VLAN ID & IP Src & IP Dst & IP Prot & TCP sport & TCP dport & Action \\
  \hline
  s2 & * & * & * & * & * & 10.1.0.1 & 10.1.0.3 & UDP & * & * & deliver as is \\
  \hline
  s2 & * & * & * & * & * & * & * & * & * & * & drop \\
  \hline

\end{tabular}
\end{center}

%%% End document
\end{document}

%%% Local Variables:
%%% mode: latex
%%% TeX-master: t
%%% End:
